% !TeX root = ../main.tex
% Add the above to each chapter to make compiling the PDF easier in some editors.

\chapter{Introduction}\label{chapter:introduction}

  \section{Motivation}
  The human brain is one of the most complex and important systems within the human body. The brain manages various functionalities like emotion, memory, attention and sensory processing, to name a few. Apart from that, the brain plays a significant role in how we process musical pieces. Independent of culture and age, music is a universal tool to influence one's emotions, concentration and cognition. Furthermore, the brain gives us the unique ability to retrieve many different features out of music, like melody extraction or the identification of an artist or a specific genre. Depending on the style of music, the brain reacts differently. However, although the brain is central to cognitive processes such as listening to music, the core mechanisms remain unknown and are currently black boxes.\par
  To enlighten these black boxes, medical advancements like \ac{fMRI} or \ac{EEG} arrived, which can measure brain activity. Especially EEG has aroused high interest since they directly measure the neural activity and can be done non-invasive with a little overhead in setup. This is why combining EEG and music has the ability to help better understand how our brain functions. By further adding means of \ac{ML}, a whole new interdisciplinary field opens up, bridging the gap between neurosciences and Computer science.\par
  There is a growing research interest in the affective influence of music on the brain. However, classifying different music genres based on \ac{EEG} signals remains underrepresented. This is exactly where this thesis aims to close the gap, by investigating whether it is possible to classify various musical genres depending on the EEG signals acquired using state-of-the-art ML techniques.\par
  The research results have high potential in the domain of medical computer science. Especially with the growing investigations on \acp{BMI}, adaptive music therapy, personalised music recommendations or neurofeedback systems may become possible. Additionally, entertainment technologies can become more user-centred, facilitating unique experiences that impact many sectors like education and mental health \par



  \section{Research goals}

  \section{Methodology}
